\chapter{Statistics for Machine Learning}
\label{ch:stats-for-ml}

Chapter 14 turned data into estimates: maximise a likelihood, or combine it with a prior.
Machine learning does the same thing at a larger scale, with its own notation and its own
habits. Open a paper and you meet $\E_{x\sim p}[\ell(f_\theta(x),y)]$, $\hat\theta$,
$\argmin$, ``validation AUC $0.91$''. Every one of those symbols is a statistical idea you
already own. This chapter is the decoder ring, plus the handful of statistical tools that ML
practice leans on daily: honest error estimates, the bias--variance tradeoff, regularisation,
classifier evaluation, and error bars on the metrics themselves.

\section{A notation decoder}
\label{sec:ml-notation}

Here is a sentence lifted from the style of a typical paper:
\[
  \hat\theta \;=\; \argmin_{\theta}\; \frac1n\sum_{i=1}^n \ell\big(f_\theta(x_i),\,y_i\big),
  \qquad (x_i,y_i)\iid p .
\]
Read aloud: ``theta-hat is the parameter value that makes the average loss on the $n$
training examples as small as possible, where the examples were drawn independently from the
same distribution $p$.'' Each piece deserves a picture.

\begin{itemize}
  \item A \textbf{loss} $\ell(\hat y, y)$ is a number that says how bad a prediction $\hat y$
  is when the truth is $y$. Squared loss $(\hat y-y)^2$ is the one from regression (Chapter 12).
  \item $f_\theta$ is a \textbf{model}: a formula with adjustable knobs $\theta$ (the
  \textbf{parameters}) that turns an input $x$ into a prediction.
  \item The hat on $\hat\theta$ means ``computed from data''. It is an estimate, and it would
  come out differently on a different sample. Same convention as $\hat p$ and $\hat\beta$ in
  earlier chapters.
  \item $\argmin_\theta g(\theta)$ is the \emph{location} of the minimum, not the minimum
  value. If $g(\theta)=(\theta-3)^2+5$ then $\min g = 5$ but $\argmin g = 3$.
\end{itemize}

Table~\ref{tab:ml-decoder} collects the symbols you will meet most often.

\begin{table}[H]
\centering
\small
\renewcommand{\arraystretch}{1.25}
\begin{tabular}{p{3.3cm} p{3.5cm} p{6.4cm}}
\toprule
\textbf{Symbol} & \textbf{Read it as} & \textbf{What it means} \\
\midrule
$x\sim p$ & ``$x$ drawn from $p$'' & $x$ is a random draw from the distribution $p$. \\
$(x_i,y_i)\iid p$ & ``i.i.d.\ from $p$'' & Independent draws, all from the same $p$. Each example tells you nothing about the next. \\
$\E_{x\sim p}[f(x)]$ & ``expected $f$ under $p$'' & The long-run average of $f(x)$ if you kept drawing $x$ from $p$. \\
$\frac1n\sum_i f(x_i)$ & ``sample average'' & The data's estimate of $\E_{x\sim p}[f(x)]$. \\
$\hat p$ or $\hat p_n$ & ``empirical distribution'' & Put weight $1/n$ on each observed point. Averages under $\hat p$ are sample averages. \\
$p(y\given x;\theta)$ & ``prob.\ of $y$ given $x$, at $\theta$'' & A model's predicted distribution for the label. The semicolon separates data from parameters. \\
$\hat\theta$, $\hat y$ & ``theta-hat'', ``y-hat'' & Quantities estimated or predicted from data. \\
$\theta^\star$ & ``theta-star'' & The best possible parameter value (usually unknowable). \\
$\argmin_\theta L(\theta)$ & ``the $\theta$ minimising $L$'' & Where the minimum occurs, not its value. \\
$\nabla_\theta L$ & ``gradient of $L$'' & Vector of slopes of $L$ in each parameter direction. \\
$\Sigma$ & ``covariance matrix'' & Variances on the diagonal, covariances off it (Section~\ref{sec:ml-cov}). \\
\bottomrule
\end{tabular}
\caption{Decoder for common ML notation. Nearly every entry is an expectation, a sample
average, or an estimate of one.}
\label{tab:ml-decoder}
\end{table}

\subsection{True distribution versus empirical distribution}

The distinction that matters most in the table is $p$ versus $\hat p$. The \textbf{true
distribution} $p$ is the process that generates data: every email that will ever arrive,
every future house sale. You never see it. The \textbf{empirical distribution} $\hat p_n$ is
the data you actually hold, viewed as a distribution that puts probability $1/n$ on each
point.

Take $n=12$ draws from a standard Normal:
\[
-1.95,\ -1.69,\ -1.41,\ -0.30,\ -0.08,\ -0.04,\ 0.06,\ 0.21,\ 0.39,\ 0.50,\ 0.59,\ 1.87.
\]
The empirical CDF climbs by $1/12$ at each point. Figure~\ref{fig:ml-ecdf} lays it over the
true CDF. The staircase tracks the smooth curve, imperfectly. With more data the steps get
finer and the fit tightens. This is the Glivenko--Cantelli theorem that justified the
bootstrap in Chapter 8; Probability From the Ground Up has the proof.

\begin{figure}[H]
\centering
\begin{tikzpicture}
\begin{axis}[
  width=11cm, height=5.6cm,
  xmin=-3, xmax=3, ymin=0, ymax=1.08,
  xlabel={$x$}, ylabel={cumulative probability $\P(X\le x)$},
  axis lines=left, legend pos=north west, legend style={font=\small},
]
\addplot[densityred, thick, domain=-3:3, samples=120] {1/(1+exp(-1.702*x))};
\addlegendentry{true CDF of $p=\Normal(0,1)$}
\addplot[sourceblue, thick, const plot] coordinates {
(-3,0) (-1.95,0.083) (-1.69,0.167) (-1.41,0.25) (-0.30,0.333) (-0.08,0.417)
(-0.04,0.5) (0.06,0.583) (0.21,0.667) (0.39,0.75) (0.50,0.833) (0.59,0.917)
(1.87,1) (3,1)};
\addlegendentry{empirical CDF of $\hat p_{12}$}
\addplot[only marks, mark=|, mark size=4pt, sourceblue] coordinates {
(-1.95,0) (-1.69,0) (-1.41,0) (-0.30,0) (-0.08,0) (-0.04,0) (0.06,0) (0.21,0) (0.39,0)
(0.50,0) (0.59,0) (1.87,0)};
\node[font=\small, anchor=west, align=left] at (axis cs:0.75,0.52) {each observation lifts\\the staircase by $1/12$};
\draw[->, gray] (axis cs:0.95,0.63) -- (axis cs:0.63,0.86);
\node[font=\small, anchor=west, align=left, sourceblue] at (axis cs:0.85,0.24) {ticks on the axis:\\the $12$ observations};
\draw[->, gray] (axis cs:1.55,0.12) -- (axis cs:1.83,0.04);
\end{axis}
\end{tikzpicture}
\caption{The empirical distribution $\hat p_{12}$ (blue staircase, one step of $1/12$ at each
observed value) is the data's stand-in for the unknown true distribution $p$ (smooth red
curve). The two agree roughly but not exactly. Training a model on data means optimising
under $\hat p$ while hoping to do well under $p$.}
\label{fig:ml-ecdf}
\end{figure}

\begin{intuitionbox}
Any expectation under $\hat p_n$ is just a sample average:
$\E_{x\sim\hat p_n}[f(x)] = \frac1n\sum_i f(x_i)$. So when a paper switches between
$\E_{x\sim p}$ and $\E_{x\sim\hat p}$, it is switching between ``the average over the whole
world'' and ``the average over my dataset''. The gap between the two is what the rest of this
chapter is about.
\end{intuitionbox}

\section{Expected risk, empirical risk, and the generalisation gap}
\label{sec:ml-risk}

\subsection{A tiny regression, by hand}

Five houses: floor area $x$ (hundreds of m$^2$) and price $y$ (hundreds of thousands).
\[
x = (1,2,3,4,5),\qquad y = (2.5,\ 3.5,\ 6.5,\ 8.0,\ 9.5).
\]
Try the model $f_\theta(x)=\theta x$ with $\theta=2$. Predictions are $2,4,6,8,10$; residuals
$y-\hat y$ are $0.5$, $-0.5$, $0.5$, $0$, $-0.5$; squared losses are $0.25$, $0.25$, $0.25$, $0$, $0.25$.
Their average is
\[
\hat R(2) = \frac{0.25+0.25+0.25+0+0.25}{5} = 0.20 .
\]
That number is the \textbf{empirical risk}: the average loss on the data you have
(Figure~\ref{fig:ml-residuals}).

\begin{figure}[H]
\centering
\begin{tikzpicture}
\begin{axis}[
  width=9.5cm, height=7.2cm,
  xmin=0, xmax=5.8, ymin=0, ymax=11.5,
  xlabel={floor area $x$ (hundreds of m$^2$)}, ylabel={price $y$ (hundreds of thousands)},
  axis lines=left, xtick={1,2,3,4,5}, ytick={0,2,4,6,8,10},
  legend pos=north west, legend style={font=\small},
]
\addplot[sourceblue, thick, domain=0:5.6] {2*x};
\addlegendentry{model $f_2(x)=2x$}
\addplot[only marks, mark=*, mark size=1.8pt, densityred] coordinates {(1,2.5) (2,3.5) (3,6.5) (4,8.0) (5,9.5)};
\addlegendentry{the five houses}
\draw[structurepurple, line width=1.8pt] (axis cs:1,2.5) -- (axis cs:1,2);
\draw[structurepurple, line width=1.8pt] (axis cs:2,3.5) -- (axis cs:2,4);
\draw[structurepurple, line width=1.8pt] (axis cs:3,6.5) -- (axis cs:3,6);
\draw[structurepurple, line width=1.8pt] (axis cs:5,9.5) -- (axis cs:5,10);
\node[font=\small, structurepurple, anchor=east] at (axis cs:0.92,2.25) {$0.25$};
\node[font=\small, structurepurple, anchor=west] at (axis cs:2.08,3.6) {$0.25$};
\node[font=\small, structurepurple, anchor=east] at (axis cs:2.92,6.25) {$0.25$};
\node[font=\small, structurepurple, anchor=west] at (axis cs:4.08,7.4) {$0$};
\node[font=\small, structurepurple, anchor=west] at (axis cs:5.08,9.6) {$0.25$};
\end{axis}
\end{tikzpicture}
\caption{Empirical risk as a picture. Each purple segment is a residual $y-\hat y$ of size
$0.5$ (house 4 sits exactly on the line), and the number beside it is its squared loss
$0.5^2=0.25$. The average of the five numbers, $1.0/5=0.20$, is $\hat R(2)$. A different
$\theta$ tilts the line and changes all five numbers.}
\label{fig:ml-residuals}
\end{figure}

\begin{definition}[Expected and empirical risk]
For a model $f$ and loss $\ell$, the \textbf{expected risk} (also called true risk or
population risk) is the average loss on a fresh example from the true distribution:
\[
R(f) = \E_{(x,y)\sim p}\big[\ell(f(x),y)\big].
\]
The \textbf{empirical risk} is the same average over the $n$ training examples:
\[
\hat R_n(f) = \frac1n\sum_{i=1}^n \ell(f(x_i),y_i) = \E_{(x,y)\sim\hat p_n}\big[\ell(f(x),y)\big].
\]
\end{definition}

For any \emph{fixed} model chosen before seeing the data, $\hat R_n(f)$ is a sample mean of
i.i.d.\ losses. So it is unbiased for $R(f)$, and its standard error is
$\SD(\ell)/\sqrt n$ (Chapter 6). The trouble starts when the model is \emph{chosen using} the
same data.

\subsection{Empirical risk minimisation}

\textbf{Empirical risk minimisation} (ERM) picks the model with the smallest training loss:
\[
\hat f = \argmin_{f\in\mathcal F} \hat R_n(f),
\]
where $\mathcal F$ is the family of models you allow (all lines, all depth-5 trees, all
networks of a given shape). For our houses and $f_\theta(x)=\theta x$, setting the derivative
of $\sum(y_i-\theta x_i)^2$ to zero gives
\[
\hat\theta = \frac{\sum x_iy_i}{\sum x_i^2} = \frac{108.5}{55} \approx 1.973,
\qquad \hat R_5(\hat\theta)\approx 0.192 .
\]
ERM found a slightly better $\theta$ than our guess of 2, and its training loss dropped from
$0.200$ to $0.192$. Least squares (Chapter 12) and maximum likelihood (Chapter 14) are both
ERM: least squares with squared loss, maximum likelihood with loss $-\log p(y\given x;\theta)$.

\begin{whybox}
Why is the training loss of $\hat f$ optimistic? Because $\hat f$ was picked to make
$\hat R_n$ small. Some of the small value is real skill and some is luck: the chosen model
happens to fit the particular noise in this sample. A fresh sample has different noise, so on
average over samples $\E[\hat R_n(\hat f)] \le \E[R(\hat f)]$. The same thing
happened with $s^2$ in Chapter 2: measuring spread around $\bar x$, a point fitted to the
data, underestimates spread around $\mu$, and the $n-1$ corrects for it.
\end{whybox}

\begin{definition}[Generalisation gap]
The \textbf{generalisation gap} of a fitted model is $R(\hat f) - \hat R_n(\hat f)$: how much
worse it does on new data than on the data it was trained on.
\end{definition}

A flexible model can drive the training loss to zero by memorising. A polynomial of degree
$n-1$ passes through all $n$ points exactly. Its training loss is $0$ and its gap is
everything. Figure~\ref{fig:ml-gap} shows the typical pattern as the model gets more
flexible.

\begin{figure}[H]
\centering
\begin{tikzpicture}
\begin{axis}[
  width=11cm, height=6cm, clip=false,
  xmin=1, xmax=12, ymin=0, ymax=1.1,
  xlabel={model flexibility (e.g.\ polynomial degree)}, ylabel={mean squared error},
  axis lines=left, legend pos=north east, legend style={font=\small},
  xtick={1,3,5,7,9,11},
]
\addplot[sourceblue, thick, domain=1:12, samples=100] {0.9*exp(-0.55*(x-1)) + 0.09*exp(-0.15*(x-1))};
\addlegendentry{training error $\hat R_n(\hat f)$}
\addplot[densityred, thick, domain=1:12, samples=100] {0.9*exp(-0.55*(x-1)) + 0.02*exp(0.3*(x-1)) + 0.09};
\addlegendentry{test error $\approx R(\hat f)$}
\draw[<->, structurepurple, thick] (axis cs:10,0.0297) -- (axis cs:10,0.394);
\node[structurepurple, anchor=west, font=\small, align=left] at (axis cs:10.1,0.21) {generalisation\\gap};
\draw[dashed, gray] (axis cs:6.2,0) -- (axis cs:6.2,0.237);
\node[gray, anchor=north, font=\small, fill=white, inner sep=1pt] at (axis cs:6.2,0.55) {best test error};
\draw[gray, ->] (axis cs:6.2,0.47) -- (axis cs:6.2,0.26);
\end{axis}
\end{tikzpicture}
\caption{The typical shape (schematic curves). Training error keeps falling as flexibility
grows; test error falls, bottoms out, then climbs. The vertical distance between them is the
generalisation gap, and it widens as the model gets more flexible. Training error alone
always says ``more flexible is better''.}
\label{fig:ml-gap}
\end{figure}

\begin{recoverybox}
Pause and think: a model scores $0.0$ training error on $1000$ examples. What, if anything,
does that tell you about its test error? (Almost nothing. It tells you the model is flexible
enough to memorise $1000$ examples.)
\end{recoverybox}

\section{Honest error estimates: splits and cross-validation}
\label{sec:ml-splits}

\subsection{Train, validation, test}

The fix for optimistic training error is the same idea as an experiment's control group: score
the model on data it has never touched. Standard practice cuts the data into three parts
(Figure~\ref{fig:ml-split}).

\begin{itemize}
  \item \textbf{Training set}: used to fit the parameters $\theta$.
  \item \textbf{Validation set}: used to choose among models and \textbf{hyperparameters}
  (settings you pick rather than fit, such as polynomial degree or regularisation strength).
  \item \textbf{Test set}: touched once, at the end, to report performance.
\end{itemize}

\begin{figure}[H]
\centering
\begin{tikzpicture}[x=1cm,y=1cm]
\fill[sourceblue!30] (0,0) rectangle (7.2,1);
\fill[quotientorange!35] (7.2,0) rectangle (9.6,1);
\fill[targetgreen!35] (9.6,0) rectangle (12,1);
\draw[thick] (0,0) rectangle (12,1);
\draw[thick] (7.2,0) -- (7.2,1); \draw[thick] (9.6,0) -- (9.6,1);
\node[font=\small] at (3.6,0.5) {training $60\%$};
\node[font=\small, align=center] at (8.4,0.5) {validation\\$20\%$};
\node[font=\small, align=center] at (10.8,0.5) {test\\$20\%$};
\node[font=\small, align=center, anchor=north] at (3.6,-0.15) {used to\\fit $\theta$};
\node[font=\small, align=center, anchor=north] at (8.4,-0.15) {used to pick\\degree, $\lambda$, \ldots};
\node[font=\small, align=center, anchor=north] at (10.8,-0.15) {used once,\\to report};
\draw[->, thick, gray] (3.6,1.15) to[bend left=25] node[above, font=\small, black] {candidate fitted models} (8.2,1.15);
\draw[->, thick, gray] (8.6,1.15) to[bend left=35] node[above, font=\small, black, pos=0.6] {one chosen model} (10.8,1.15);
\end{tikzpicture}
\caption{A three-way split. Each part answers one question. The test score is honest only if
nothing about the model was chosen by looking at the test set.}
\label{fig:ml-split}
\end{figure}

Why does the validation set not suffice for reporting? Because you used it to choose. If you
try $50$ hyperparameter settings and report the best validation score, you have run ERM over
those $50$ settings on the validation set, and the winner's score is optimistic for the same
reason training error is. This is the multiple-comparisons problem of Chapter 11 again.

\subsection{$k$-fold cross-validation}

With $500$ examples, holding out $20\%$ for validation wastes data and gives a noisy score.
\textbf{$k$-fold cross-validation} reuses every example for both jobs. Cut the data into $k$
equal \textbf{folds}. For each fold in turn, train on the other $k-1$ folds and score on the
held-out fold. Average the $k$ scores (Figure~\ref{fig:ml-kfold}).

\begin{figure}[H]
\centering
\begin{tikzpicture}[x=1cm,y=0.62cm]
\foreach \r in {1,...,5} {
  \foreach \c in {1,...,5} {
    \ifnum\r=\c
      \fill[quotientorange!45] (1.6*\c,-\r) rectangle (1.6*\c+1.6,-\r+0.8);
    \else
      \fill[sourceblue!25] (1.6*\c,-\r) rectangle (1.6*\c+1.6,-\r+0.8);
    \fi
    \draw (1.6*\c,-\r) rectangle (1.6*\c+1.6,-\r+0.8);
  }
  \node[anchor=east, font=\small] at (1.5,-\r+0.4) {round \r};
}
\foreach \c in {1,...,5} { \node[font=\small] at (1.6*\c+0.8,0.4) {fold \c}; }
\node[anchor=west, font=\small] at (9.8,-0.6) {$0.21$};
\node[anchor=west, font=\small] at (9.8,-1.6) {$0.25$};
\node[anchor=west, font=\small] at (9.8,-2.6) {$0.19$};
\node[anchor=west, font=\small] at (9.8,-3.6) {$0.23$};
\node[anchor=west, font=\small] at (9.8,-4.6) {$0.22$};
\node[anchor=west, font=\small, align=left] at (9.8,0.5) {held-out\\error};
\draw (9.8,-5.1) -- (11,-5.1);
\node[anchor=west, font=\small, align=left] at (9.8,-5.75) {CV estimate:\\mean $0.22$};
\fill[quotientorange!45] (1.6,-6.2) rectangle (1.95,-5.85); \draw (1.6,-6.2) rectangle (1.95,-5.85);
\node[anchor=west, font=\small] at (2.0,-6.03) {held out and scored};
\fill[sourceblue!25] (5.8,-6.2) rectangle (6.15,-5.85); \draw (5.8,-6.2) rectangle (6.15,-5.85);
\node[anchor=west, font=\small] at (6.2,-6.03) {used for training};
\end{tikzpicture}
\caption{Five-fold cross-validation. In each round one fold (orange) is held out and scored
while the model is trained on the other four (blue). Each example is held out exactly once,
and the five held-out errors are averaged into one CV estimate.}
\label{fig:ml-kfold}
\end{figure}

With the five fold errors $0.21, 0.25, 0.19, 0.23, 0.22$, the CV estimate is their mean,
$0.22$. Their standard deviation is $0.022$. The naive standard error $0.022/\sqrt5=0.010$ is
too small: the five training sets overlap by $75\%$, so the fold scores are positively
correlated, and the averaging buys less than independence would. Treat the fold-to-fold spread
as a rough guide to stability, not as a valid confidence interval.

Common choices are $k=5$ or $k=10$. $k=n$ is \textbf{leave-one-out} CV: nearly unbiased but
slow and high-variance. In all cases the final model is refit on all the data after CV has
chosen its settings.

\subsection{Data leakage}

\textbf{Data leakage} is any route by which information from the held-out data reaches the
model during training. It makes held-out scores optimistic, sometimes wildly so. Classic
leaks:

\begin{itemize}
  \item Standardising features (subtracting the mean, dividing by the SD) using \emph{all}
  the data, then splitting. The training set has now seen the test set's mean.
  \item Selecting the ``$20$ most predictive features'' on all the data, then cross-validating
  a model on those features.
  \item Random splits of time-series data, so the model trains on Thursday and is tested on
  Wednesday.
  \item Duplicate or near-duplicate rows (the same patient twice, the same image resized)
  landing on both sides of the split.
  \item A feature that is a consequence of the label, such as ``refund issued'' when
  predicting fraud.
\end{itemize}

\begin{figure}[H]
\centering
\begin{tikzpicture}[
  box/.style={draw, rounded corners, minimum height=0.8cm, align=center, font=\small, inner sep=3pt},
  arr/.style={->, thick}]
\node[box, fill=trapred!12] (a1) at (0,0) {all data};
\node[box, fill=trapred!12] (a2) at (3.0,0) {fit scaler /\\select features};
\node[box, fill=trapred!12] (a3) at (6.2,0) {split};
\node[box, fill=trapred!12] (a4) at (9.8,0) {train, test\\(optimistic)};
\draw[arr] (a1) -- (a2); \draw[arr] (a2) -- (a3); \draw[arr] (a3) -- (a4);
\node[anchor=east, trapred, font=\small\bfseries] at (-0.9,0) {leaky};
\node[box, fill=targetgreen!12] (b1) at (0,-2) {all data};
\node[box, fill=targetgreen!12] (b2) at (3.0,-2) {split};
\node[box, fill=targetgreen!12] (b3) at (6.2,-2) {fit scaler /\\select on train};
\node[box, fill=targetgreen!12] (b4) at (9.8,-2) {apply to test,\\score (honest)};
\draw[arr] (b1) -- (b2); \draw[arr] (b2) -- (b3); \draw[arr] (b3) -- (b4);
\node[anchor=east, targetgreen, font=\small\bfseries] at (-0.9,-2) {clean};
\end{tikzpicture}
\caption{The order of operations decides whether the test score is honest. Anything that
learns from data (scaling, imputation, feature selection) belongs inside the training side of
the split, and inside each CV fold.}
\label{fig:ml-leak}
\end{figure}

\begin{trapbox}
Feature selection before cross-validation is the most damaging leak because it looks
innocent. With $5000$ pure-noise features and $50$ examples, picking the $20$ features most
correlated with the label on all the data and then cross-validating can report accuracy far
above chance on labels that are coin flips. The selection already saw every fold.
\end{trapbox}

\section{The bias--variance tradeoff for models}
\label{sec:ml-biasvar}

Chapter 14 split the error of an estimator into bias and variance:
$\MSE = \Bias^2 + \Var$. Bias is how far the estimator's average lands from the target;
variance is how much it scatters from sample to sample. A fitted model is an estimator too.
It estimates the true regression function $f(x)=\E[y\given x]$, the average $y$ at input $x$,
and the same split applies at every $x$.

\subsection{Many fits, one picture}

Take the true curve $f(x)=\sin(2\pi x)$ on $[0,1]$. At $15$ evenly spaced inputs
$x=\tfrac1{30},\tfrac3{30},\dots,\tfrac{29}{30}$, draw noisy values $y=f(x)+\eps$ with noise
SD $\sigma=0.3$, fit a model, and repeat on fresh datasets. Each dataset gives a different
fitted curve. Figure~\ref{fig:ml-manyfits} fits a straight line and a degree-9 polynomial to
the same six datasets.

\begin{figure}[H]
\centering
\begin{tikzpicture}
\begin{axis}[
  name=left, width=6.6cm, height=5.4cm,
  xmin=0, xmax=1, ymin=-2.5, ymax=2.5,
  xlabel={$x$}, ylabel={$y$},
  title={degree 1: high bias, low variance}, title style={font=\small},
  axis lines=left, clip=true, xtick={0,0.5,1},
]
\addplot[sourceblue!80, thin] coordinates {(0.000,0.88) (0.025,0.84) (0.050,0.79) (0.075,0.75) (0.100,0.70) (0.125,0.66) (0.150,0.62) (0.175,0.57) (0.200,0.53) (0.225,0.48) (0.250,0.44) (0.275,0.40) (0.300,0.35) (0.325,0.31) (0.350,0.26) (0.375,0.22) (0.400,0.18) (0.425,0.13) (0.450,0.09) (0.475,0.04) (0.500,0.00) (0.525,-0.04) (0.550,-0.09) (0.575,-0.13) (0.600,-0.18) (0.625,-0.22) (0.650,-0.26) (0.675,-0.31) (0.700,-0.35) (0.725,-0.40) (0.750,-0.44) (0.775,-0.48) (0.800,-0.53) (0.825,-0.57) (0.850,-0.61) (0.875,-0.66) (0.900,-0.70) (0.925,-0.75) (0.950,-0.79) (0.975,-0.83) (1.000,-0.88)};
\addplot[sourceblue!80, thin] coordinates {(0.000,0.86) (0.025,0.81) (0.050,0.77) (0.075,0.73) (0.100,0.69) (0.125,0.64) (0.150,0.60) (0.175,0.56) (0.200,0.51) (0.225,0.47) (0.250,0.43) (0.275,0.38) (0.300,0.34) (0.325,0.30) (0.350,0.26) (0.375,0.21) (0.400,0.17) (0.425,0.13) (0.450,0.08) (0.475,0.04) (0.500,-0.00) (0.525,-0.05) (0.550,-0.09) (0.575,-0.13) (0.600,-0.17) (0.625,-0.22) (0.650,-0.26) (0.675,-0.30) (0.700,-0.35) (0.725,-0.39) (0.750,-0.43) (0.775,-0.48) (0.800,-0.52) (0.825,-0.56) (0.850,-0.60) (0.875,-0.65) (0.900,-0.69) (0.925,-0.73) (0.950,-0.78) (0.975,-0.82) (1.000,-0.86)};
\addplot[sourceblue!80, thin] coordinates {(0.000,1.01) (0.025,0.96) (0.050,0.92) (0.075,0.87) (0.100,0.82) (0.125,0.77) (0.150,0.73) (0.175,0.68) (0.200,0.63) (0.225,0.58) (0.250,0.54) (0.275,0.49) (0.300,0.44) (0.325,0.39) (0.350,0.35) (0.375,0.30) (0.400,0.25) (0.425,0.21) (0.450,0.16) (0.475,0.11) (0.500,0.06) (0.525,0.02) (0.550,-0.03) (0.575,-0.08) (0.600,-0.13) (0.625,-0.17) (0.650,-0.22) (0.675,-0.27) (0.700,-0.32) (0.725,-0.36) (0.750,-0.41) (0.775,-0.46) (0.800,-0.51) (0.825,-0.55) (0.850,-0.60) (0.875,-0.65) (0.900,-0.70) (0.925,-0.74) (0.950,-0.79) (0.975,-0.84) (1.000,-0.89)};
\addplot[sourceblue!80, thin] coordinates {(0.000,1.45) (0.025,1.39) (0.050,1.33) (0.075,1.27) (0.100,1.21) (0.125,1.14) (0.150,1.08) (0.175,1.02) (0.200,0.96) (0.225,0.90) (0.250,0.84) (0.275,0.77) (0.300,0.71) (0.325,0.65) (0.350,0.59) (0.375,0.53) (0.400,0.46) (0.425,0.40) (0.450,0.34) (0.475,0.28) (0.500,0.22) (0.525,0.15) (0.550,0.09) (0.575,0.03) (0.600,-0.03) (0.625,-0.09) (0.650,-0.15) (0.675,-0.22) (0.700,-0.28) (0.725,-0.34) (0.750,-0.40) (0.775,-0.46) (0.800,-0.53) (0.825,-0.59) (0.850,-0.65) (0.875,-0.71) (0.900,-0.77) (0.925,-0.83) (0.950,-0.90) (0.975,-0.96) (1.000,-1.02)};
\addplot[sourceblue!80, thin] coordinates {(0.000,0.62) (0.025,0.58) (0.050,0.55) (0.075,0.51) (0.100,0.47) (0.125,0.44) (0.150,0.40) (0.175,0.37) (0.200,0.33) (0.225,0.29) (0.250,0.26) (0.275,0.22) (0.300,0.19) (0.325,0.15) (0.350,0.11) (0.375,0.08) (0.400,0.04) (0.425,0.01) (0.450,-0.03) (0.475,-0.06) (0.500,-0.10) (0.525,-0.14) (0.550,-0.17) (0.575,-0.21) (0.600,-0.24) (0.625,-0.28) (0.650,-0.32) (0.675,-0.35) (0.700,-0.39) (0.725,-0.42) (0.750,-0.46) (0.775,-0.50) (0.800,-0.53) (0.825,-0.57) (0.850,-0.60) (0.875,-0.64) (0.900,-0.67) (0.925,-0.71) (0.950,-0.75) (0.975,-0.78) (1.000,-0.82)};
\addplot[sourceblue!80, thin] coordinates {(0.000,1.07) (0.025,1.02) (0.050,0.96) (0.075,0.91) (0.100,0.86) (0.125,0.80) (0.150,0.75) (0.175,0.70) (0.200,0.64) (0.225,0.59) (0.250,0.54) (0.275,0.48) (0.300,0.43) (0.325,0.38) (0.350,0.32) (0.375,0.27) (0.400,0.22) (0.425,0.16) (0.450,0.11) (0.475,0.06) (0.500,0.00) (0.525,-0.05) (0.550,-0.10) (0.575,-0.16) (0.600,-0.21) (0.625,-0.26) (0.650,-0.32) (0.675,-0.37) (0.700,-0.42) (0.725,-0.48) (0.750,-0.53) (0.775,-0.58) (0.800,-0.64) (0.825,-0.69) (0.850,-0.74) (0.875,-0.80) (0.900,-0.85) (0.925,-0.90) (0.950,-0.96) (0.975,-1.01) (1.000,-1.06)};
\addplot[black, very thick, dashed, domain=0:1, samples=80] {sin(360*x)};
\node[font=\small, anchor=north east, fill=white, inner sep=1pt] at (axis cs:0.98,2.4) {dashed: truth $\sin(2\pi x)$};
\end{axis}
\begin{axis}[
  at={(left.east)}, anchor=west, xshift=1.2cm,
  width=6.6cm, height=5.4cm,
  xmin=0, xmax=1, ymin=-2.5, ymax=2.5,
  xlabel={$x$},
  title={degree 9: low bias, high variance}, title style={font=\small},
  axis lines=left, clip=true, xtick={0,0.5,1},
]
\addplot[trapred!85, thin] coordinates {(0.000,-3.26) (0.025,-0.48) (0.050,0.50) (0.075,0.66) (0.100,0.56) (0.125,0.47) (0.150,0.49) (0.175,0.62) (0.200,0.82) (0.225,1.01) (0.250,1.17) (0.275,1.24) (0.300,1.23) (0.325,1.14) (0.350,0.97) (0.375,0.77) (0.400,0.55) (0.425,0.34) (0.450,0.15) (0.475,-0.01) (0.500,-0.14) (0.525,-0.25) (0.550,-0.33) (0.575,-0.42) (0.600,-0.50) (0.625,-0.60) (0.650,-0.69) (0.675,-0.79) (0.700,-0.87) (0.725,-0.92) (0.750,-0.94) (0.775,-0.91) (0.800,-0.85) (0.825,-0.78) (0.850,-0.71) (0.875,-0.67) (0.900,-0.67) (0.925,-0.66) (0.950,-0.50) (0.975,0.15) (1.000,1.85)};
\addplot[trapred!85, thin] coordinates {(0.000,1.40) (0.025,0.17) (0.050,-0.05) (0.075,0.16) (0.100,0.49) (0.125,0.77) (0.150,0.95) (0.175,1.02) (0.200,1.00) (0.225,0.93) (0.250,0.84) (0.275,0.76) (0.300,0.69) (0.325,0.63) (0.350,0.59) (0.375,0.55) (0.400,0.50) (0.425,0.43) (0.450,0.34) (0.475,0.22) (0.500,0.08) (0.525,-0.07) (0.550,-0.23) (0.575,-0.37) (0.600,-0.50) (0.625,-0.60) (0.650,-0.67) (0.675,-0.71) (0.700,-0.72) (0.725,-0.71) (0.750,-0.71) (0.775,-0.71) (0.800,-0.73) (0.825,-0.78) (0.850,-0.83) (0.875,-0.86) (0.900,-0.85) (0.925,-0.73) (0.950,-0.46) (0.975,0.01) (1.000,0.69)};
\addplot[trapred!85, thin] coordinates {(0.000,-1.39) (0.025,-0.40) (0.050,0.28) (0.075,0.72) (0.100,0.98) (0.125,1.12) (0.150,1.18) (0.175,1.19) (0.200,1.18) (0.225,1.15) (0.250,1.11) (0.275,1.06) (0.300,0.99) (0.325,0.89) (0.350,0.78) (0.375,0.65) (0.400,0.51) (0.425,0.36) (0.450,0.21) (0.475,0.07) (0.500,-0.05) (0.525,-0.16) (0.550,-0.25) (0.575,-0.33) (0.600,-0.41) (0.625,-0.49) (0.650,-0.59) (0.675,-0.70) (0.700,-0.84) (0.725,-0.98) (0.750,-1.12) (0.775,-1.22) (0.800,-1.26) (0.825,-1.18) (0.850,-0.96) (0.875,-0.60) (0.900,-0.15) (0.925,0.25) (0.950,0.33) (0.975,-0.41) (1.000,-2.73)};
\addplot[trapred!85, thin] coordinates {(0.000,-6.01) (0.025,-0.48) (0.050,1.39) (0.075,1.56) (0.100,1.19) (0.125,0.83) (0.150,0.71) (0.175,0.84) (0.200,1.13) (0.225,1.47) (0.250,1.76) (0.275,1.93) (0.300,1.95) (0.325,1.82) (0.350,1.58) (0.375,1.26) (0.400,0.92) (0.425,0.59) (0.450,0.31) (0.475,0.09) (0.500,-0.07) (0.525,-0.18) (0.550,-0.27) (0.575,-0.35) (0.600,-0.45) (0.625,-0.58) (0.650,-0.73) (0.675,-0.89) (0.700,-1.05) (0.725,-1.16) (0.750,-1.19) (0.775,-1.14) (0.800,-0.99) (0.825,-0.76) (0.850,-0.52) (0.875,-0.32) (0.900,-0.24) (0.925,-0.30) (0.950,-0.41) (0.975,-0.25) (1.000,0.86)};
\addplot[trapred!85, thin] coordinates {(0.000,-2.48) (0.025,-0.44) (0.050,0.28) (0.075,0.42) (0.100,0.37) (0.125,0.32) (0.150,0.33) (0.175,0.42) (0.200,0.54) (0.225,0.67) (0.250,0.78) (0.275,0.84) (0.300,0.86) (0.325,0.83) (0.350,0.77) (0.375,0.69) (0.400,0.59) (0.425,0.48) (0.450,0.36) (0.475,0.21) (0.500,0.05) (0.525,-0.14) (0.550,-0.36) (0.575,-0.59) (0.600,-0.83) (0.625,-1.04) (0.650,-1.22) (0.675,-1.33) (0.700,-1.36) (0.725,-1.29) (0.750,-1.13) (0.775,-0.90) (0.800,-0.66) (0.825,-0.47) (0.850,-0.39) (0.875,-0.46) (0.900,-0.67) (0.925,-0.85) (0.950,-0.62) (0.975,0.77) (1.000,4.59)};
\addplot[trapred!85, thin] coordinates {(0.000,-1.86) (0.025,0.06) (0.050,0.62) (0.075,0.63) (0.100,0.54) (0.125,0.51) (0.150,0.60) (0.175,0.78) (0.200,0.98) (0.225,1.16) (0.250,1.27) (0.275,1.30) (0.300,1.24) (0.325,1.12) (0.350,0.95) (0.375,0.76) (0.400,0.58) (0.425,0.41) (0.450,0.27) (0.475,0.15) (0.500,0.04) (0.525,-0.06) (0.550,-0.18) (0.575,-0.33) (0.600,-0.49) (0.625,-0.68) (0.650,-0.87) (0.675,-1.05) (0.700,-1.19) (0.725,-1.27) (0.750,-1.28) (0.775,-1.22) (0.800,-1.10) (0.825,-0.96) (0.850,-0.86) (0.875,-0.81) (0.900,-0.84) (0.925,-0.84) (0.950,-0.58) (0.975,0.48) (1.000,3.29)};
\addplot[black, very thick, dashed, domain=0:1, samples=80] {sin(360*x)};
\end{axis}
\end{tikzpicture}
\caption{The same six datasets of $15$ noisy points, fitted two ways; dashed black is the
truth. Left: the lines agree with each other but all miss the wave (bias). Right: the
polynomials follow the wave on average but disagree wildly with each other, especially near
the edges, where some shoot off the plot (variance).}
\label{fig:ml-manyfits}
\end{figure}

Repeat this on $20{,}000$ simulated datasets. At each $x$, the average of the fitted curves
gives the bias (its distance from the truth) and their spread gives the variance. Averaging
both over $x\in[0,1]$, with $\sigma^2=0.09$:

\begin{center}
\begin{tabular}{lccc}
\toprule
model & bias$^2$ & variance & bias$^2$ + variance + $\sigma^2$ \\
\midrule
degree 1 & $0.200$ & $0.012$ & $0.302$ \\
degree 3 & $0.005$ & $0.025$ & $0.120$ \\
degree 9 & $0.000$ & $0.209$ & $0.299$ \\
\bottomrule
\end{tabular}
\end{center}

Degree 3 wins. It has enough flexibility to bend once up and once down, and not enough to chase
noise. Degrees 1 and 9 score about the same, for opposite reasons.

\begin{theorem}[Bias--variance decomposition]
Let $y=f(x)+\eps$ with $\E[\eps]=0$, $\Var(\eps)=\sigma^2$, and let $\hat f$ be fitted on a
random training set independent of the new noise $\eps$. At a fixed input $x$,
\[
\E\big[(y-\hat f(x))^2\big]
= \underbrace{\sigma^2}_{\text{noise}}
+ \underbrace{\big(\E[\hat f(x)]-f(x)\big)^2}_{\text{bias}^2}
+ \underbrace{\Var\big(\hat f(x)\big)}_{\text{variance}} ,
\]
where the expectation is over both the training set and the new noise.
\end{theorem}

\begin{whybox}
Write $y-\hat f = \eps + (f-\E\hat f) + (\E\hat f-\hat f)$. The three pieces are: fresh noise,
a constant offset, and the fit's wobble around its own average. Square and take expectations.
The cross terms vanish: $\eps$ is independent of the fit and has mean zero, and
$\E\hat f-\hat f$ has mean zero by construction. What remains is the sum of the three squared
pieces. No model can beat the noise term $\sigma^2$; it is the floor.
\end{whybox}

Plotted against flexibility, bias$^2$ falls, variance rises, and their sum with the noise floor
traces the U of Figure~\ref{fig:ml-biasvar}.

\begin{figure}[H]
\centering
\begin{tikzpicture}
\begin{axis}[
  width=11.5cm, height=6cm,
  xmin=1, xmax=12, ymin=0, ymax=1.1,
  xlabel={model flexibility}, ylabel={expected squared error},
  axis lines=left, legend style={at={(0.5,0.98)}, anchor=north, font=\small},
  xtick={1,3,5,7,9,11},
]
\addplot[densityred, very thick, domain=1:12, samples=100] {0.9*exp(-0.55*(x-1)) + 0.02*exp(0.3*(x-1)) + 0.09};
\addlegendentry{test error (sum)}
\addplot[sourceblue, thick, domain=1:12, samples=100] {0.9*exp(-0.55*(x-1))};
\addlegendentry{bias$^2$}
\addplot[quotientorange, thick, domain=1:12, samples=100] {0.02*exp(0.3*(x-1))};
\addlegendentry{variance}
\addplot[gray, dashed, thick, domain=1:12] {0.09};
\addlegendentry{noise $\sigma^2$}
\addplot[only marks, mark=*, densityred] coordinates {(6.2,0.237)};
\node[anchor=south, font=\small] at (axis cs:6.2,0.27) {lowest test error};
\node[anchor=north west, font=\small, align=left] at (axis cs:2.0,1.05) {underfitting\\(high bias)};
\node[anchor=north east, font=\small, align=right] at (axis cs:11.9,1.05) {overfitting\\(high variance)};
\end{axis}
\end{tikzpicture}
\caption{The U-shaped test error is the sum of a falling bias$^2$ curve, a rising variance
curve, and a constant noise floor. Cross-validation is a way of locating the bottom of the U
from data.}
\label{fig:ml-biasvar}
\end{figure}

\begin{connectionbox}
Very large neural networks sometimes show \emph{double descent}: test error rises near the
point where the model can just barely interpolate the data, then falls again as the model
grows further. It does not contradict the decomposition above, which always holds. It says the
variance curve need not keep rising forever once training finds the smoothest of the many
interpolating solutions. Goodfellow et al.\ Ch.~5 discusses capacity in the classical picture.
\end{connectionbox}

\section{Regularisation: ridge and lasso}
\label{sec:ml-reg}

If variance is the problem, add bias on purpose. \textbf{Regularisation} adds a penalty to the
training loss that discourages large coefficients:
\begin{align*}
\text{ridge:}\quad \hat\beta^{\text{ridge}} &= \argmin_\beta\; \textstyle\sum_i (y_i - x_i^\top\beta)^2 + \lambda\sum_j \beta_j^2,\\
\text{lasso:}\quad \hat\beta^{\text{lasso}} &= \argmin_\beta\; \textstyle\sum_i (y_i - x_i^\top\beta)^2 + \lambda\sum_j |\beta_j|.
\end{align*}
Here $\lambda\ge0$ is a hyperparameter: the price per unit of coefficient size. $\lambda=0$
gives ordinary least squares. Large $\lambda$ pushes every coefficient toward $0$. The
features should be standardised first so that ``size'' means the same thing for each.

\subsection{Shrinkage by hand}

In the simplest case, where the features are uncorrelated and scaled to unit length
(an orthonormal design), both penalties act on each least-squares coefficient
$\hat\beta_j$ separately:
\[
\hat\beta_j^{\text{ridge}} = \frac{\hat\beta_j}{1+\lambda},
\qquad
\hat\beta_j^{\text{lasso}} = \operatorname{sign}(\hat\beta_j)\,\max\!\big(|\hat\beta_j| - \tfrac{\lambda}{2},\,0\big).
\]
Ridge shrinks every coefficient by the same factor. Lasso subtracts a fixed amount and stops at
zero (\textbf{soft thresholding}). Take least-squares coefficients $(3.0,\ 1.0,\ -0.4)$ and
$\lambda=1$:

\begin{center}
\begin{tabular}{lccc}
\toprule
& $\beta_1$ & $\beta_2$ & $\beta_3$ \\
\midrule
least squares & $3.0$ & $1.0$ & $-0.4$ \\
ridge, $\lambda=1$ (divide by 2) & $1.5$ & $0.5$ & $-0.2$ \\
lasso, $\lambda=1$ (subtract 0.5, floor at 0) & $2.5$ & $0.5$ & $0$ \\
\bottomrule
\end{tabular}
\end{center}

Lasso removed the small coefficient entirely. That is \textbf{variable selection} for free.
Figure~\ref{fig:ml-shrink} draws both rules as functions.

\begin{figure}[H]
\centering
\begin{tikzpicture}
\begin{axis}[
  width=8.5cm, height=6.5cm,
  xmin=-3, xmax=3.3, ymin=-3, ymax=3, clip=false,
  xlabel={least-squares coefficient $\hat\beta_j$}, ylabel={regularised coefficient},
  axis lines=middle, legend pos=south east, legend style={font=\small},
  xlabel style={at={(axis description cs:0.5,-0.08)}, anchor=north},
  ylabel style={at={(axis description cs:-0.08,0.5)}, anchor=south, rotate=90},
  xtick={-2,-1,1,2}, ytick={-2,-1,1,2},
  xticklabels={$-2$,,,$2$}, yticklabels={$-2$,,,$2$},
  ticklabel style={font=\small, fill=white, inner sep=1pt},
]
\addplot[gray, dashed, domain=-3:3] {x};
\addlegendentry{no penalty}
\addplot[sourceblue, thick, domain=-3:3] {x/2};
\addlegendentry{ridge}
\addplot[densityred, thick] coordinates {(-3,-2.5) (-0.5,0) (0.5,0) (3,2.5)};
\addlegendentry{lasso}
\addplot[only marks, mark=*, mark size=1.8pt, sourceblue, forget plot] coordinates {(3,1.5) (1,0.5) (-0.4,-0.2)};
\addplot[only marks, mark=*, mark size=1.8pt, densityred, forget plot] coordinates {(3,2.5) (-0.4,0)};
\addplot[only marks, mark=o, mark size=2.6pt, densityred, forget plot] coordinates {(1,0.5)};
\node[font=\small, anchor=west, sourceblue] at (axis cs:3.05,1.5) {$1.5$};
\node[font=\small, anchor=west, densityred] at (axis cs:3.05,2.5) {$2.5$};
\node[font=\small, anchor=east, align=right] at (axis cs:-0.45,1.0) {lasso dead zone:\\$|\hat\beta_j|\le0.5\mapsto0$};
\draw[->, gray] (axis cs:-0.6,0.55) -- (axis cs:-0.2,0.08);
\end{axis}
\end{tikzpicture}
\caption{Shrinkage rules for $\lambda=1$ with an orthonormal design; dots mark the worked
example $\hat\beta=(3.0,\ 1.0,\ -0.4)$. Ridge (blue) halves every coefficient, so it never
reaches zero. Lasso (red) moves every coefficient $0.5$ toward zero and sets anything in
$[-0.5,0.5]$ exactly to zero, which is why $-0.4$ becomes $0$.}
\label{fig:ml-shrink}
\end{figure}

\subsection{Why lasso gives exact zeros: the geometry}

Each penalised problem is equivalent to a constrained one: minimise the squared error subject
to $\sum\beta_j^2\le t$ (ridge) or $\sum|\beta_j|\le t$ (lasso), for some budget $t$ that
shrinks as $\lambda$ grows. In two dimensions the ridge budget is a disc and the lasso budget
is a diamond. The squared-error loss has elliptical contours centred on the least-squares
solution. The answer is the first point where a growing ellipse touches the budget region.

\begin{figure}[H]
\centering
\begin{tikzpicture}[scale=1.35]
\begin{scope}
\clip (-1.5,-1.4) rectangle (3.5,3.4);
\draw[->, gray] (-1.4,0) -- (3.4,0) node[anchor=north east, black, font=\small] {$\beta_1$};
\draw[->, gray] (0,-1.3) -- (0,3.3) node[anchor=north west, black, font=\small] {$\beta_2$};
\fill[sourceblue!25] (0,0) circle (1);
\draw[sourceblue, thick] (0,0) circle (1);
\foreach \s in {0.35,0.7,1} {
  \draw[densityred, rotate around={-30:(2,1.4)}] (2,1.4) ellipse ({0.819*\s} and {1.638*\s});
}
\fill (2,1.4) circle (1.3pt) node[anchor=south west, font=\small] {$\hat\beta^{\text{LS}}$};
\fill[structurepurple] (0.957,0.289) circle (1.6pt);
\draw[structurepurple, ->] (0.35,0.62) -- (0.9,0.32);
\node[structurepurple, anchor=east, font=\small] at (0.4,0.65) {ridge solution};
\end{scope}
\node[font=\small, align=center] at (1,-1.8) {ridge budget: disc $\beta_1^2+\beta_2^2\le1$\\both coefficients shrink, neither is zero};
\begin{scope}[xshift=5.4cm]
\clip (-1.5,-1.4) rectangle (3.5,3.4);
\draw[->, gray] (-1.4,0) -- (3.4,0) node[anchor=north east, black, font=\small] {$\beta_1$};
\draw[->, gray] (0,-1.3) -- (0,3.3) node[anchor=north west, black, font=\small] {$\beta_2$};
\fill[targetgreen!25] (1,0) -- (0,1) -- (-1,0) -- (0,-1) -- cycle;
\draw[targetgreen!70!black, thick] (1,0) -- (0,1) -- (-1,0) -- (0,-1) -- cycle;
\foreach \s in {0.35,0.7,1} {
  \draw[densityred, rotate around={-30:(2,1.4)}] (2,1.4) ellipse ({0.872*\s} and {1.744*\s});
}
\fill (2,1.4) circle (1.3pt) node[anchor=south west, font=\small] {$\hat\beta^{\text{LS}}$};
\fill[structurepurple] (1,0) circle (1.6pt);
\node[structurepurple, anchor=north west, font=\small, align=left] at (0.85,-0.3) {lasso solution\\$(1,0)$: $\beta_2=0$};
\end{scope}
\node[font=\small, align=center] at (6.4,-1.8) {lasso budget: diamond $|\beta_1|+|\beta_2|\le1$\\the contour hits a corner on the $\beta_1$ axis};
\end{tikzpicture}
\caption{Constraint geometry. Red ellipses are contours of the squared error around the
least-squares fit $\hat\beta^{\text{LS}}=(2,1.4)$. Left: the smallest contour touching the
disc meets it at $(0.96,0.29)$. Right: it meets the diamond at the corner $(1,0)$. Corners sit
on the axes, so lasso solutions often have exact zeros.}
\label{fig:ml-constraint}
\end{figure}

\begin{intuitionbox}
A disc has no corners, so the touching point can be anywhere on its rim and will almost never
land exactly on an axis. A diamond is all corners and flat edges, and in high dimensions the
corners and low-dimensional edges stick out. Expanding ellipses tend to hit them first. Each
corner is a point where some coefficients are zero.
\end{intuitionbox}

\begin{connectionbox}
Regularisation is a prior in disguise (Chapter 14). With Normal errors, ridge is the MAP
estimate under independent Normal priors $\beta_j\sim\Normal(0,\tau^2)$, and $\lambda$ is the
ratio $\sigma^2/\tau^2$. Lasso is the MAP estimate under Laplace priors, with density
proportional to $e^{-|\beta_j|/b}$: two exponential tails glued back to back, with a sharp peak
at zero. That peak plays the role of the diamond's corner.
Weight decay in neural networks is ridge applied to every weight.
\end{connectionbox}

\section{Evaluating classifiers}
\label{sec:ml-classify}

\subsection{The confusion matrix}

A spam filter scores $1000$ emails. $100$ are spam (the \textbf{positive} class) and $900$
are not. It flags $110$ emails: $80$ spam and $30$ legitimate. Every prediction falls in one
of four cells (Figure~\ref{fig:ml-confusion}).

\begin{figure}[H]
\centering
\begin{tikzpicture}[x=2.6cm, y=1.3cm]
\fill[targetgreen!25] (0,1) rectangle (1,2);
\fill[trapred!18] (1,1) rectangle (2,2);
\fill[trapred!18] (0,0) rectangle (1,1);
\fill[targetgreen!25] (1,0) rectangle (2,1);
\draw[thick] (0,0) rectangle (2,2); \draw[thick] (1,0) -- (1,2); \draw[thick] (0,1) -- (2,1);
\node[align=center, font=\small] at (0.5,1.5) {TP $=80$\\true positive};
\node[align=center, font=\small] at (1.5,1.5) {FP $=30$\\false positive};
\node[align=center, font=\small] at (0.5,0.5) {FN $=20$\\false negative};
\node[align=center, font=\small] at (1.5,0.5) {TN $=870$\\true negative};
\node[font=\small] at (0.5,2.25) {actually spam};
\node[font=\small] at (1.5,2.25) {actually not};
\node[font=\small, anchor=east] at (-0.05,1.5) {flagged};
\node[font=\small, anchor=east] at (-0.05,0.5) {not flagged};
\node[font=\small, anchor=west] at (2.05,1.5) {$110$ flagged};
\node[font=\small, anchor=west] at (2.05,0.5) {$890$ passed};
\node[font=\small] at (0.5,-0.3) {$100$ spam};
\node[font=\small] at (1.5,-0.3) {$900$ not spam};
\end{tikzpicture}
\caption{Confusion matrix for the spam filter. Green cells are correct, red cells are errors.
Every classification metric is a ratio of these four counts.}
\label{fig:ml-confusion}
\end{figure}

From the four counts:
\begin{align*}
\text{accuracy} &= \frac{TP+TN}{\text{all}} = \frac{80+870}{1000} = 0.95,\\
\text{precision} &= \frac{TP}{TP+FP} = \frac{80}{110} \approx 0.727 \quad\text{(of the flagged, how many are spam)},\\
\text{recall} = \text{TPR} &= \frac{TP}{TP+FN} = \frac{80}{100} = 0.80 \quad\text{(of the spam, how much was caught)},\\
\text{FPR} &= \frac{FP}{FP+TN} = \frac{30}{900} \approx 0.033 \quad\text{(of the good mail, how much was flagged)}.
\end{align*}
Recall is also called \textbf{sensitivity} or \textbf{true positive rate} (TPR).
$1-\text{FPR}=0.967$ is the \textbf{specificity}. The \textbf{F1 score} combines precision $P$
and recall $R$ through their harmonic mean:
\[
F_1 = \frac{2PR}{P+R} = \frac{2(0.727)(0.80)}{0.727+0.80} \approx 0.762 .
\]
The harmonic mean is dragged toward the smaller of the two, so $F_1$ is high only when both
are.

\begin{trapbox}
Accuracy $0.95$ sounds good until you notice that a filter that flags nothing scores
$900/1000=0.90$. When classes are imbalanced, always compare against the ``predict the
majority class'' baseline, and prefer precision and recall.
\end{trapbox}

\subsection{Scores, thresholds, and the ROC curve}

Most classifiers output a \textbf{score} (often a predicted probability) and flag an example
when the score exceeds a \textbf{threshold}. Moving the threshold trades one error for the
other (Figure~\ref{fig:ml-scores}).

\begin{figure}[H]
\centering
\begin{tikzpicture}
\begin{axis}[
  width=11.5cm, height=6cm,
  xmin=-3.5, xmax=4.5, ymin=0, ymax=0.56,
  xlabel={classifier score}, ylabel={density},
  axis lines=left, legend columns=2, legend style={font=\small, at={(0.5,1.04)}, anchor=south},
  ytick=\empty,
]
\addplot[area legend, fill=trapred!25, draw=none, domain=0.8:4.5, samples=60] {0.39894*exp(-x^2/2)} \closedcycle;
\addlegendentry{false positives (negatives flagged)}
\addplot[area legend, fill=quotientorange!35, draw=none, domain=-3.5:0.8, samples=60] {0.39894*exp(-(x-1.2)^2/2)} \closedcycle;
\addlegendentry{false negatives (positives missed)}
\addplot[sourceblue, thick, domain=-3.5:4.5, samples=120] {0.39894*exp(-x^2/2)};
\addlegendentry{scores of negatives (not spam)}
\addplot[densityred, thick, domain=-3.5:4.5, samples=120] {0.39894*exp(-(x-1.2)^2/2)};
\addlegendentry{scores of positives (spam)}
\draw[very thick, black] (axis cs:0.8,0) -- (axis cs:0.8,0.47);
\node[anchor=south, font=\small] at (axis cs:0.8,0.47) {threshold $0.8$};
\draw[->, thick] (axis cs:0.9,0.445) -- (axis cs:2.2,0.445) node[anchor=west, font=\small] {flagged};
\draw[->, thick] (axis cs:0.7,0.445) -- (axis cs:-0.6,0.445) node[anchor=east, font=\small] {not flagged};
\end{axis}
\end{tikzpicture}
\caption{Score distributions for negatives (blue) and positives (red) overlap. Everything right
of the threshold is flagged. Sliding it right shrinks the false positives (red shading) and
grows the false negatives (orange shading). No threshold removes both.}
\label{fig:ml-scores}
\end{figure}

The \textbf{ROC curve} (receiver operating characteristic) plots TPR against FPR for every
possible threshold. Here are $10$ test examples sorted by score, with true labels ($+$ = positive):

\begin{center}
\begin{tabular}{l*{10}{c}}
\toprule
score & $0.95$ & $0.90$ & $0.80$ & $0.70$ & $0.60$ & $0.55$ & $0.45$ & $0.40$ & $0.30$ & $0.20$\\
label & $+$ & $+$ & $-$ & $+$ & $+$ & $-$ & $-$ & $+$ & $-$ & $-$\\
\bottomrule
\end{tabular}
\end{center}

There are $5$ positives and $5$ negatives. Start with the threshold above every score: nothing
flagged, $(\text{FPR},\text{TPR})=(0,0)$. Lower the threshold past one example at a time. Each
positive moves the point up by $1/5$; each negative moves it right by $1/5$. The path is
\[
\begin{gathered}
(0,0)\to(0,0.2)\to(0,0.4)\to(0.2,0.4)\to(0.2,0.6)\to(0.2,0.8)\\
\to(0.4,0.8)\to(0.6,0.8)\to(0.6,1)\to(0.8,1)\to(1,1).
\end{gathered}
\]

\begin{figure}[H]
\centering
\begin{tikzpicture}
\begin{axis}[
  width=7.5cm, height=7.5cm,
  xmin=0, xmax=1, ymin=0, ymax=1.02,
  xlabel={false positive rate}, ylabel={true positive rate},
  axis lines=left, xtick={0,0.2,0.4,0.6,0.8,1}, ytick={0,0.2,0.4,0.6,0.8,1},
  grid=major, grid style={gray!20},
]
\addplot[fill=histfill!45, draw=none] coordinates {
(0,0) (0,0.4) (0.2,0.4) (0.2,0.8) (0.6,0.8) (0.6,1) (1,1) (1,0)} -- cycle;
\addplot[distblue, very thick, mark=*, mark size=1.6pt] coordinates {
(0,0) (0,0.2) (0,0.4) (0.2,0.4) (0.2,0.6) (0.2,0.8) (0.4,0.8) (0.6,0.8) (0.6,1) (0.8,1) (1,1)};
\addplot[gray, dashed, domain=0:1] {x};
\node[font=\small, fill=white, inner sep=1pt] at (axis cs:0.62,0.42) {AUC $=0.80$};
\node[font=\small, gray, rotate=45] at (axis cs:0.8,0.7) {random};
\end{axis}
\end{tikzpicture}
\caption{ROC curve for the $10$ scored examples. Each up-step is a positive, each right-step a
negative, taken in score order. The shaded area is $0.80$. The dashed diagonal is a classifier
that guesses at random.}
\label{fig:ml-roc}
\end{figure}

\begin{definition}[AUC]
The \textbf{AUC} is the area under the ROC curve. Equivalently, it is the probability that a
randomly chosen positive gets a higher score than a randomly chosen negative (ties count
one half).
\end{definition}

Check the equivalence on our data. There are $5\times5=25$ positive--negative pairs. Count the
pairs in which the positive outscores the negative: the positive at $0.95$ beats all $5$
negatives, $0.90$ beats $5$, $0.70$ beats $4$, $0.60$ beats $4$, and $0.40$ beats $2$. Total
$20$, and $20/25 = 0.80$, matching the area. The area interpretation is a picture; the pairs
interpretation is the one to remember. It is the Mann--Whitney $U$ statistic of Chapter 10
divided by the number of pairs.

\begin{recoverybox}
Pause and think: what AUC does a classifier get if it outputs the same score for every
example? (One half. Every pair is a tie.) What if you flip the sign of every score of a model
with AUC $0.80$? ($0.20$.)
\end{recoverybox}

\begin{historybox}
ROC curves come from radar. During the Second World War, operators watched noisy screens and
had to call each blip ``aircraft'' or ``noise''. A jumpy operator caught every plane and raised
many false alarms; a cautious one missed planes. Engineers wanted to compare the
\emph{receivers} (the radar sets) without mixing in each operator's caution. Their answer was to
plot hit rate against false-alarm rate across all levels of caution. In 1954 Peterson, Birdsall
and Fox turned this into signal detection theory, and psychologists Green and Swets used it in
the 1960s to separate what a listener can hear from how willing they are to say ``yes''.
Hanley and McNeil (1982) showed medical researchers that the area under the curve equals the
probability of ranking a random sick patient above a random healthy one. That is the AUC every
ML dashboard now reports.
\end{historybox}

\subsection{Precision--recall curves and base rates}

ROC curves ignore how common positives are: TPR is computed within positives, FPR within
negatives. Precision mixes the two, so it depends on the \textbf{base rate} (the fraction of
examples that are positive). The \textbf{precision--recall (PR) curve} plots precision against
recall across thresholds. For our $10$ examples, flagging the top $k$ for $k=1,\dots,10$ gives
recall $=(\text{positives flagged})/5$ and precision $=(\text{positives flagged})/k$:

\begin{figure}[H]
\centering
\begin{tikzpicture}
\begin{axis}[
  width=8.5cm, height=6.8cm,
  xmin=0, xmax=1.05, ymin=0, ymax=1.12, clip=false,
  xlabel={recall}, ylabel={precision},
  axis lines=left, grid=major, grid style={gray!20},
  xtick={0,0.2,0.4,0.6,0.8,1}, ytick={0,0.2,0.4,0.6,0.8,1},
]
\addplot[densityred, very thick, mark=*, mark size=1.6pt] coordinates {
(0.2,1.0) (0.4,1.0) (0.4,0.667) (0.6,0.75) (0.8,0.8) (0.8,0.667) (0.8,0.571) (1.0,0.625) (1.0,0.556) (1.0,0.5)};
\addplot[gray, dashed, domain=0:1.05] {0.5};
\node[gray, font=\small, anchor=north west] at (axis cs:0.03,0.45) {base rate $0.5$ (random)};
\node[font=\footnotesize, anchor=south, inner sep=2.5pt] at (axis cs:0.2,1.0) {$k=1$};
\node[font=\footnotesize, anchor=south, inner sep=2.5pt] at (axis cs:0.4,1.0) {$k=2$};
\node[font=\footnotesize, anchor=east, inner sep=2.5pt] at (axis cs:0.4,0.667) {$k=3$};
\node[font=\footnotesize, anchor=south, inner sep=2.5pt] at (axis cs:0.6,0.75) {$k=4$};
\node[font=\footnotesize, anchor=south, inner sep=2.5pt] at (axis cs:0.8,0.8) {$k=5$};
\node[font=\footnotesize, anchor=east, inner sep=2.5pt] at (axis cs:0.8,0.667) {$k=6$};
\node[font=\footnotesize, anchor=east, inner sep=2.5pt] at (axis cs:0.8,0.571) {$k=7$};
\node[font=\footnotesize, anchor=west, align=left] at (axis cs:1.03,0.565) {$k=8,9,10$\\(top to bottom)};
\end{axis}
\end{tikzpicture}
\caption{Precision--recall curve for the same $10$ examples; dot $k$ means ``flag the top $k$
scores''. It zigzags: each positive flagged moves right and up, each negative flagged drops
straight down. A random classifier's precision equals the base rate (dashed), so the floor of
this plot moves with how rare positives are.}
\label{fig:ml-pr}
\end{figure}

Base rates can turn a good classifier into a disappointing one. Keep the spam filter's
$\text{TPR}=0.8$ and $\text{FPR}=1/30$, and deploy it where only $1\%$ of $100{,}000$ emails
are spam. It catches $0.8\times1000 = 800$ spam and flags $99{,}000/30 = 3300$ good emails.
Precision is
\[
\frac{800}{800+3300} \approx 0.195 .
\]
Four out of five flags are now wrong, with the same ROC point. In general
\[
\text{precision} = \frac{\text{TPR}\cdot\pi}{\text{TPR}\cdot\pi + \text{FPR}\cdot(1-\pi)},
\]
where $\pi$ is the base rate. This is Bayes' rule (Chapter 14) with the classifier as the test.

\begin{figure}[H]
\centering
\begin{tikzpicture}
\begin{axis}[
  width=10.5cm, height=5.8cm,
  xmode=log, xmin=0.001, xmax=0.5, ymin=0, ymax=1,
  xlabel={base rate $\pi$ (fraction positive, log scale)}, ylabel={precision},
  axis lines=left, grid=major, grid style={gray!20}, clip=false,
  xtick={0.001,0.01,0.1,0.5}, xticklabels={$0.001$,$0.01$,$0.1$,$0.5$},
]
\addplot[structurepurple, very thick, domain=0.001:0.5, samples=120] {0.8*x/(0.8*x + (1-x)/30)};
\addplot[only marks, mark=*, densityred] coordinates {(0.1,0.727) (0.01,0.195)};
\node[font=\small, anchor=east] at (axis cs:0.09,0.75) {test set ($\pi=0.1$): $0.727$};
\node[font=\small, anchor=west] at (axis cs:0.011,0.17) {deployment ($\pi=0.01$): $0.195$};
\end{axis}
\end{tikzpicture}
\caption{Precision of a fixed classifier ($\text{TPR}=0.8$, $\text{FPR}=1/30$) as the base rate
changes. Same model, same ROC point, wildly different precision. Report the base rate with
every precision number.}
\label{fig:ml-baserate}
\end{figure}

Use ROC/AUC when you care about ranking quality independent of class balance. Use PR curves
when positives are rare and the cost of a false alarm matters, as in fraud, search, and
medical screening.

\subsection{Calibration}

A weather forecaster who says ``$70\%$ chance of rain'' is \textbf{calibrated} if it rains on
about $70\%$ of such days. A classifier is calibrated if, among examples it scores near $0.7$,
about $70\%$ are positive. AUC cannot see calibration: squaring every score leaves the ranking,
and hence the AUC, unchanged, while wrecking the probabilities.

To check, sort $1000$ test predictions into ten bins by predicted probability and compare each
bin's mean prediction with the fraction of positives in it. The result is a
\textbf{reliability diagram} (Figure~\ref{fig:ml-calib}).

\begin{figure}[H]
\centering
\begin{tikzpicture}
\begin{axis}[
  width=7.8cm, height=7.8cm,
  xmin=0, xmax=1, ymin=0, ymax=1,
  xlabel={mean predicted probability in bin}, ylabel={observed fraction positive},
  axis lines=left, grid=major, grid style={gray!20}, clip=false,
  legend style={font=\small, at={(1.04,0.0)}, anchor=south west},
  xtick={0,0.2,0.4,0.6,0.8,1}, ytick={0,0.2,0.4,0.6,0.8,1},
]
\addplot[gray, dashed, domain=0:1] {x};
\addlegendentry{perfect calibration}
\addplot[ybar, area legend, bar width=0.07, fill=histfill!70, draw=distblue] coordinates {
(0.05,0.12) (0.15,0.20) (0.25,0.31) (0.35,0.39) (0.45,0.47) (0.55,0.53) (0.65,0.59) (0.75,0.66) (0.85,0.72) (0.95,0.80)};
\addlegendentry{model}
\addplot[densityred, very thick] coordinates {(0.05,0.05) (0.05,0.12)};
\addlegendentry{gap (calibration error)}
\addplot[densityred, very thick, forget plot] coordinates {(0.15,0.15) (0.15,0.20)};
\addplot[densityred, very thick, forget plot] coordinates {(0.25,0.25) (0.25,0.31)};
\addplot[densityred, very thick, forget plot] coordinates {(0.35,0.35) (0.35,0.39)};
\addplot[densityred, very thick, forget plot] coordinates {(0.45,0.45) (0.45,0.47)};
\addplot[densityred, very thick, forget plot] coordinates {(0.55,0.55) (0.55,0.53)};
\addplot[densityred, very thick, forget plot] coordinates {(0.65,0.65) (0.65,0.59)};
\addplot[densityred, very thick, forget plot] coordinates {(0.75,0.75) (0.75,0.66)};
\addplot[densityred, very thick, forget plot] coordinates {(0.85,0.85) (0.85,0.72)};
\addplot[densityred, very thick, forget plot] coordinates {(0.95,0.95) (0.95,0.80)};
\node[font=\small, anchor=west, align=left] at (axis cs:1.04,0.86) {top bin:\\predicted $0.95$,\\observed $0.80$};
\draw[->, gray] (axis cs:1.05,0.8) -- (axis cs:0.99,0.8);
\node[font=\small, anchor=west, align=left] at (axis cs:0.02,0.7) {bottom bin:\\predicted $0.05$,\\observed $0.12$};
\draw[->, gray] (axis cs:0.06,0.58) -- (axis cs:0.05,0.14);
\end{axis}
\end{tikzpicture}
\caption{Reliability diagram for an overconfident model. Each red segment is one bin's gap
between the bar (what happened) and the diagonal (what was predicted). Predictions of $0.95$
come true only $80\%$ of the time, and predictions of $0.05$ come true $12\%$ of the time: the
model's probabilities are too extreme at both ends.}
\label{fig:ml-calib}
\end{figure}

A one-number summary is the \textbf{expected calibration error} (ECE): the bin-size-weighted
average of $|\text{mean prediction}-\text{observed fraction}|$. With bin sizes
$300,120,80,60,50,50,60,80,100,100$ for the ten bins above, the gaps are
$0.07,0.05,0.06,0.04,0.02,0.02,0.06,0.09,0.13,0.15$, and
\[
\text{ECE} = \frac{300(0.07)+120(0.05)+\cdots+100(0.15)}{1000} = \frac{75}{1000} = 0.075 .
\]
Overconfidence of this shape is common in large neural networks. A cheap fix is to refit a
one-parameter rescaling of the scores on the validation set (temperature scaling, a close cousin
of logistic regression from Chapter 13).

\section{Uncertainty in ML metrics}
\label{sec:ml-uncertainty}

A test-set metric is a statistic computed from a sample. It has a standard error. A
leaderboard that reports ``$87.0\%$ vs $86.4\%$'' without one is reporting noise as often as
signal.

\subsection{A confidence interval for accuracy}

Each test example is a Bernoulli trial: right ($1$) or wrong ($0$). If the test examples are
i.i.d., the number correct is Binomial, and accuracy is a sample proportion. Chapter 7's
interval applies directly. With $435$ correct out of $n=500$, $\hat a = 0.87$:
\[
\SE = \sqrt{\frac{\hat a(1-\hat a)}{n}} = \sqrt{\frac{0.87\times0.13}{500}} \approx 0.0150,
\qquad
0.87 \pm 1.96\times0.0150 = [0.841,\ 0.899].
\]
The true accuracy could plausibly be anywhere from $84\%$ to $90\%$. To pin accuracy near
$0.9$ to $\pm1$ point at $95\%$ confidence you need
$n\approx 1.96^2(0.9)(0.1)/0.01^2\approx 3460$ test examples. For accuracies near $0$ or $1$
or small $n$, use the Wilson interval instead (second line of output below).

\begin{lstlisting}[language=Python]
import numpy as np
from scipy import stats
k, n = 435, 500
a = k / n; se = np.sqrt(a * (1 - a) / n)
print(a - 1.96 * se, a + 1.96 * se)          # Wald: 0.841, 0.899
print(stats.binomtest(k, n).proportion_ci(method="wilson"))  # 0.838, 0.897
\end{lstlisting}

\subsection{A bootstrap interval for AUC}

AUC has no simple standard-error formula worth memorising. The bootstrap (Chapter 8) handles
it: resample the test set with replacement, recompute AUC, repeat, and read off percentiles.
Resample positives and negatives separately so every replicate keeps the same class balance.

On a test set of $60$ positives and $60$ negatives, a model scores AUC $=0.769$. Two thousand
stratified bootstrap replicates give the histogram in Figure~\ref{fig:ml-bootauc}. The middle
$95\%$ runs from $0.680$ to $0.852$, and the bootstrap SE is $0.043$.

\begin{figure}[H]
\centering
\begin{tikzpicture}
\begin{axis}[
  width=11cm, height=5.6cm,
  xmin=0.56, xmax=0.95, ymin=0, ymax=470, clip=false,
  xlabel={bootstrap AUC}, ylabel={count (of 2000)},
  axis lines=left, xtick={0.60,0.65,0.70,0.75,0.80,0.85,0.90},
  xticklabel style={/pgf/number format/fixed, /pgf/number format/precision=2},
]
\addplot[ybar interval, fill=bootorange!55, draw=bootorange!80!black] coordinates {
(0.60,2) (0.62,2) (0.64,16) (0.66,31) (0.68,63) (0.70,133) (0.72,231) (0.74,304)
(0.76,354) (0.78,368) (0.80,260) (0.82,141) (0.84,66) (0.86,23) (0.88,6) (0.90,0)};
\draw[densityred, very thick, dashed] (axis cs:0.680,0) -- (axis cs:0.680,465);
\draw[densityred, very thick, dashed] (axis cs:0.852,0) -- (axis cs:0.852,465);
\draw[black, thick] (axis cs:0.769,0) -- (axis cs:0.769,385);
\node[font=\small, anchor=east, densityred, align=right] at (axis cs:0.678,350) {$2.5$th percentile\\$0.680$};
\node[font=\small, anchor=west, densityred, align=left] at (axis cs:0.854,350) {$97.5$th percentile\\$0.852$};
\node[font=\small, anchor=south, inner sep=2pt] at (axis cs:0.769,385) {observed AUC $0.769$};
\draw[<->, densityred, thick] (axis cs:0.680,465) -- (axis cs:0.852,465);
\node[font=\small, anchor=south, densityred] at (axis cs:0.766,467) {middle $95\%$ of bootstrap AUCs: width $0.17$};
\end{axis}
\end{tikzpicture}
\caption{Bootstrap distribution of AUC on a $120$-example test set. The $95\%$ percentile
interval is nearly $0.17$ wide. Small test sets give AUCs that are much less precise than their
three decimal places suggest.}
\label{fig:ml-bootauc}
\end{figure}

\begin{lstlisting}[language=Python]
import numpy as np
rng = np.random.default_rng(0)
def auc(p, n):                      # P(pos score > neg score), ties = 1/2
    return np.mean(p[:, None] > n) + 0.5 * np.mean(p[:, None] == n)
def boot_auc(pos, neg, B=2000):
    reps = [auc(rng.choice(pos, len(pos)), rng.choice(neg, len(neg)))
            for _ in range(B)]
    return np.percentile(reps, [2.5, 97.5])
\end{lstlisting}

\subsection{Comparing two models on the same test set}

Model A gets $435/500$ right and model B gets $420/500$. Is A better, or did it get lucky? Two
separate confidence intervals overlap heavily, but that comparison throws away information:
both models were scored on the \emph{same} examples. This is paired data, like the
before/after measurements of Chapter 9. Only the examples where the models disagree carry
evidence about which is better.

\begin{center}
\begin{tabular}{lcc}
\toprule
& B right & B wrong \\
\midrule
A right & $395$ & $b=40$ \\
A wrong & $c=25$ & $40$ \\
\bottomrule
\end{tabular}
\end{center}

Check the margins: A is right on $395+40=435$, B on $395+25=420$. If the models were equally
good, each of the $b+c=65$ disagreements would be a fair coin flip between ``A wins'' and
``B wins''. \textbf{McNemar's test} asks whether $40$ versus $25$ is a surprising split for $65$
fair flips:
\[
\chi^2 = \frac{(b-c)^2}{b+c} = \frac{(40-25)^2}{65} = \frac{225}{65} \approx 3.46,
\qquad p \approx 0.063 \ \ (\chi^2 \text{ with 1 df}).
\]
The exact binomial version (count $\le 25$ or $\ge 40$ out of $65$ under $p=1/2$) gives
$p\approx0.082$. The three-point accuracy gap is suggestive but not convincing on $500$
examples (Figure~\ref{fig:ml-mcnemar}).

\begin{figure}[H]
\centering
\begin{tikzpicture}
\begin{axis}[
  width=11.5cm, height=5.6cm,
  xmin=14.5, xmax=50.5, ymin=0, ymax=0.125, clip=false,
  xlabel={disagreements won by model A, out of $b+c=65$},
  ylabel={probability if A and B are equal},
  ylabel style={font=\small},
  axis lines=left, xtick={15,20,25,30,35,40,45,50},
  ytick={0,0.05,0.1},
  yticklabel style={/pgf/number format/fixed, /pgf/number format/precision=2},
]
\addplot[ybar, bar width=0.75, fill=histfill!60, draw=distblue!70] coordinates {
(26,0.0272) (27,0.0393) (28,0.0533) (29,0.0680) (30,0.0816) (31,0.0921) (32,0.0978) (33,0.0978)
(34,0.0921) (35,0.0816) (36,0.0680) (37,0.0533) (38,0.0393) (39,0.0272)};
\addplot[ybar, bar width=0.75, fill=densityred!55, draw=densityred] coordinates {
(15,0.00002) (16,0.00004) (17,0.0001) (18,0.0001) (19,0.0003) (20,0.0008) (21,0.0016) (22,0.0033) (23,0.0062) (24,0.0108) (25,0.0177)
(40,0.0177) (41,0.0108) (42,0.0062) (43,0.0033) (44,0.0016) (45,0.0008) (46,0.0003) (47,0.0001) (48,0.0001) (49,0.00004) (50,0.00002)};
\draw[black, thick] (axis cs:40,0) -- (axis cs:40,0.06);
\node[font=\small, anchor=south, align=center] at (axis cs:40,0.06) {observed:\\A won $40$};
\node[font=\small, anchor=south, align=center] at (axis cs:25,0.06) {mirror image:\\B won $40$};
\draw[black, thick] (axis cs:25,0) -- (axis cs:25,0.06);
\node[font=\small, densityred, anchor=west, align=left] at (axis cs:41.2,0.045) {red: as lopsided or\\more, $p=0.082$};
\node[font=\small, anchor=south] at (axis cs:32.5,0.1) {expected split $32.5$};
\end{axis}
\end{tikzpicture}
\caption{McNemar's test as a coin-flip picture. If the models were equally good, the number of
the $65$ disagreements won by A would follow this Binomial$(65,\tfrac12)$ distribution,
centred on $32.5$. The red bars are splits at least as lopsided as the observed $40$--$25$;
together they have probability $0.082$, the exact $p$-value. Such a split is not rare enough
to be convincing.}
\label{fig:ml-mcnemar}
\end{figure}

\begin{lstlisting}[language=Python]
from scipy import stats
b, c = 40, 25                         # discordant pairs
chi2 = (b - c)**2 / (b + c)
print(chi2, stats.chi2.sf(chi2, 1))   # 3.46, 0.063
print(stats.binomtest(c, b + c, 0.5).pvalue)   # exact: 0.082
\end{lstlisting}

For metrics that are not per-example averages (AUC, F1), use a \textbf{paired bootstrap}:
resample test-example indices with replacement, compute the metric for \emph{both} models on
the same resample, record the difference, and look at the percentile interval of the
differences. If it excludes zero, the gap is unlikely to be resampling noise.

\begin{trapbox}
Comparing a new model against a baseline across $20$ random seeds and $10$ benchmarks, then
reporting the one where it wins, is the garden of forking paths from Chapter 11. Decide the
comparison in advance, and correct for the number of comparisons you actually ran.
\end{trapbox}

\section{Stochastic gradients are Monte Carlo estimates}
\label{sec:ml-sgd}

Gradient descent minimises the training loss by repeatedly stepping downhill:
$\theta \leftarrow \theta - \eta\,\nabla_\theta \hat R_n(\theta)$, where the \textbf{learning
rate} $\eta$ is the step size. The gradient of an average is the average of the gradients:
\[
\nabla_\theta \hat R_n(\theta) = \frac1n\sum_{i=1}^n \nabla_\theta\,\ell_i(\theta),
\qquad \ell_i(\theta) = \ell\big(f_\theta(x_i),y_i\big).
\]
With $n$ in the millions, computing all $n$ terms per step is wasteful. \textbf{Stochastic
gradient descent} (SGD) instead averages over a random \textbf{minibatch} of $B$ examples. That
is exactly a sample mean estimating a population mean, where the ``population'' is the $n$
per-example gradients. An average of random draws used to estimate an expectation is called a
\textbf{Monte Carlo estimate}; the bootstrap of Chapter 8 was one too.

\subsection{By hand: four points}

Fit $f_\theta(x)=\theta x$ with squared loss to $x=(1,2,3,4)$, $y=(2,5,5,9)$, at the current
value $\theta=1$. The per-example gradient is $-2x_i(y_i-\theta x_i)$:
\[
g = \big(-2(1)(1),\ -2(2)(3),\ -2(3)(2),\ -2(4)(5)\big) = (-2,\ -12,\ -12,\ -40),
\qquad \bar g = -16.5 .
\]
The six possible minibatches of size $B=2$ give gradient estimates
\[
\tfrac{-2-12}{2}=-7,\ \ \tfrac{-2-12}{2}=-7,\ \ \tfrac{-2-40}{2}=-21,\ \
\tfrac{-12-12}{2}=-12,\ \ \tfrac{-12-40}{2}=-26,\ \ \tfrac{-12-40}{2}=-26.
\]
Their average is $(-7-7-21-12-26-26)/6 = -99/6 = -16.5$, the full gradient. So the minibatch
gradient is \textbf{unbiased}. It is noisy: single examples ($B=1$) give estimates with
variance $200.75$; pairs give variance $66.9$. In general the variance falls like $1/B$, so the
standard error falls like $1/\sqrt B$, the same $\sigma/\sqrt n$ law as Chapter 6. (Here the
drop is a bit faster than $1/B$ because we sample without replacement from only four points.)
Figure~\ref{fig:ml-minibatch} lays out every possible estimate.

\begin{figure}[H]
\centering
\begin{tikzpicture}[x=0.18cm, y=1cm]
\draw[->, gray] (-43,0) -- (1.5,0);
\foreach \v in {-40,-30,-20,-10,0} { \draw[gray] (\v,0) -- (\v,-0.1) node[anchor=north, black, font=\small] {$\v$}; }
\node[font=\small, anchor=north] at (-20,-0.45) {minibatch gradient estimate};
\draw[densityred, thick, dashed] (-16.5,-0.1) -- (-16.5,2.55);
\node[densityred, font=\small, anchor=south] at (-16.5,2.55) {full gradient $\bar g=-16.5$};
\node[font=\small, anchor=east, align=right] at (-44,1.65) {$B=1$: $4$ estimates,\\variance $200.75$};
\foreach \v/\h in {-2/0, -12/0, -12/1, -40/0} { \fill[sourceblue] (\v,1.6+0.32*\h) circle (3pt); }
\draw[sourceblue, |-|] (-40,1.3) -- (-2,1.3);
\node[font=\small, anchor=east, align=right] at (-44,0.5) {$B=2$: $6$ estimates,\\variance $66.9$};
\foreach \v/\h in {-7/0, -7/1, -21/0, -12/0, -26/0, -26/1} { \fill[structurepurple] (\v,0.45+0.32*\h) circle (3pt); }
\draw[structurepurple, |-|] (-26,0.25) -- (-7,0.25);
\end{tikzpicture}
\caption{Every minibatch gradient for the four-point example; stacked dots are equal
estimates. Both rows balance on the full gradient $-16.5$ (dashed), which is what unbiased
means. Doubling the batch size pulls the estimates in toward it: the range shrinks from
$[-40,-2]$ to $[-26,-7]$.}
\label{fig:ml-minibatch}
\end{figure}

\begin{happeningbox}
Read the true objective as $R(\theta)=\E_{(x,y)\sim p}[\ell]$. Then each fresh minibatch is a
Monte Carlo estimate of $\nabla R(\theta)$ itself, not just of the training-set gradient. In one
pass over fresh data, SGD is directly descending the expected risk. That is why papers write
the gradient as $\E_{(x,y)\sim p}[\nabla_\theta\ell]$ and then approximate it with a batch
average in the next line.
\end{happeningbox}

\begin{figure}[H]
\centering
\begin{tikzpicture}[x=1.75cm, y=1.75cm]
\begin{scope}
\clip (-1.8,-0.8) rectangle (2.4,3.2);
\foreach \c/\a/\b in {0.1/0.177/0.478, 0.5/0.396/1.070, 1.5/0.686/1.853, 3.5/1.049/2.830, 7/1.483/4.002} {
  \draw[gray!60, rotate around={63.08:(1.204,1.952)}] (1.204,1.952) ellipse ({\a} and {\b});
}
\draw[sourceblue, very thick, -{Stealth[length=2mm]}] plot coordinates {
(-1.50,-0.50) (-0.17,1.73) (0.25,2.22) (0.43,2.29) (0.55,2.27) (0.64,2.23) (0.72,2.20)
(0.78,2.17) (0.84,2.14) (0.89,2.11) (0.93,2.09) (0.97,2.07) (1.00,2.06) (1.03,2.04)
(1.05,2.03) (1.07,2.02) (1.09,2.01) (1.10,2.00) (1.12,2.00) (1.13,1.99) (1.14,1.98)};
\draw[densityred, thick] plot coordinates {
(-1.50,-0.50) (0.11,2.74) (0.61,2.50) (0.84,2.67) (0.46,1.87) (0.56,1.83) (0.72,2.46)
(0.45,1.74) (0.75,2.16) (0.74,2.16) (0.87,1.97) (0.92,2.17) (0.74,1.75) (1.11,1.84)
(1.25,1.98) (0.72,1.63) (1.18,2.40) (1.18,2.40) (0.70,1.20) (1.16,2.44) (0.95,1.99)};
\fill (1.204,1.952) circle (1.5pt);
\fill (-1.5,-0.5) circle (1.5pt);
\end{scope}
\draw[gray] (-1.8,-0.8) rectangle (2.4,3.2);
\node[font=\small, anchor=north west] at (-1.5,-0.5) {start};
\draw[black, thin] (1.24,1.952) -- (1.45,1.952);
\node[font=\small, anchor=west, fill=white, inner sep=1pt] at (1.45,1.952) {minimum};
\node[font=\small, anchor=west, sourceblue, fill=white, inner sep=1pt] at (-1.75,3.0) {full-batch gradient descent};
\node[font=\small, anchor=west, densityred, fill=white, inner sep=1pt] at (-1.75,2.7) {SGD with $B=1$};
\node[font=\small] at (0.3,-1.0) {intercept $\theta_0$};
\node[font=\small, rotate=90] at (-2.0,1.2) {slope $\theta_1$};
\end{tikzpicture}
\caption{Gradient descent on a two-parameter linear regression ($200$ points), plotted every
second step. Grey ellipses are contours of the training loss. Full-batch steps (blue) glide
smoothly to the minimum. Single-example steps (red) point downhill only on average, so the path
jitters around the minimum instead of settling. Larger batches or a shrinking learning rate
reduce the jitter.}
\label{fig:ml-sgd}
\end{figure}

\section{Covariance matrices, the multivariate Normal, and PCA}
\label{sec:ml-cov}

ML inputs are vectors. A house is (area, age, distance to centre); an image is a million
pixels. The spread of a vector needs more than one number.

\subsection{The covariance matrix}

Six points in two dimensions:
\[
(2,1),\ (3,3),\ (4,3),\ (5,5),\ (6,4),\ (4,2).
\]
The mean is $(4,3)$. The deviations from the mean are
$(-2,-2),(-1,0),(0,0),(1,2),(2,1),(0,-1)$. Using the $n-1=5$ divisor from Chapter 2:
\[
s_{11} = \tfrac{4+1+0+1+4+0}{5} = 2,\qquad
s_{22} = \tfrac{4+0+0+4+1+1}{5} = 2,\qquad
s_{12} = \tfrac{4+0+0+2+2+0}{5} = 1.6 .
\]
Stack them into the \textbf{sample covariance matrix}
\[
S = \begin{pmatrix} 2 & 1.6 \\ 1.6 & 2 \end{pmatrix}.
\]
The diagonal holds each coordinate's variance. The off-diagonal holds the covariance of
Chapter 4; dividing by the SDs gives the correlation $1.6/\sqrt{2\cdot2}=0.8$. In general, for
data vectors $x_i\in\R^d$ with mean $\bar x$,
\[
S = \frac{1}{n-1}\sum_{i=1}^n (x_i-\bar x)(x_i-\bar x)^\top ,
\]
a $d\times d$ symmetric matrix. Its population version is
$\Sigma = \E[(x-\mu)(x-\mu)^\top]$.

\subsection{The multivariate Normal}

The \textbf{multivariate Normal} $\Normal(\mu,\Sigma)$ is the bell curve in $d$ dimensions. Its
density is
\[
p(x) = \frac{1}{(2\pi)^{d/2}\,|\Sigma|^{1/2}}
\exp\!\Big(-\tfrac12 (x-\mu)^\top\Sigma^{-1}(x-\mu)\Big).
\]
The exponent is a squared distance that measures each direction in units of that direction's
spread. Its contours are ellipses centred at $\mu$, with axes along the eigenvectors of
$\Sigma$ and half-lengths proportional to the square roots of the eigenvalues. When $d=1$ it
reduces to the familiar $\Normal(\mu,\sigma^2)$ density. Bishop Ch.~2 is the standard reference.

\subsection{PCA: the directions of greatest spread}

An \textbf{eigenvector} of $S$ is a direction $v$ that $S$ only stretches, never turns:
$Sv=\lambda v$, and the stretch factor $\lambda$ (no relation to the ridge penalty) is its
\textbf{eigenvalue}. For a covariance
matrix, the eigenvalue is the variance of the data measured along $v$, so the eigenvectors are
the directions along which the data are most and least spread out. For our $S$, the eigenvalues are $2+1.6=3.6$ and $2-1.6=0.4$, with eigenvectors
$(1,1)/\sqrt2$ and $(1,-1)/\sqrt2$. Check one: $S(1,1)^\top = (3.6,3.6)^\top = 3.6\,(1,1)^\top$.

\textbf{Principal component analysis} (PCA) uses these directions as new axes. The first
principal component points along $(1,1)/\sqrt2$, and the variance of the data projected onto it
is $3.6$. The total variance is $2+2=4$, so the first component captures $3.6/4=90\%$ of it.
Dropping the second component compresses each point from two numbers to one while losing only
$10\%$ of the variance.

\begin{figure}[H]
\centering
\begin{tikzpicture}[scale=0.9]
\draw[step=1, gray!20] (0.5,-0.5) grid (7.5,6.5);
\draw[->, gray] (0.5,0) -- (7.6,0) node[anchor=north east, black, font=\small] {$x_1$};
\draw[->, gray] (1,-0.5) -- (1,6.6) node[anchor=north west, black, font=\small] {$x_2$};
\foreach \x in {2,...,6} { \node[font=\footnotesize, anchor=north] at (\x,0) {\x}; }
\foreach \y in {1,...,6} { \node[font=\footnotesize, anchor=east] at (1,\y) {\y}; }
\draw[densityred, thick, rotate around={45:(4,3)}] (4,3) ellipse ({2*1.897} and {2*0.632});
\foreach \p in {(2,1),(3,3),(4,3),(5,5),(6,4),(4,2)} { \fill[sourceblue] \p circle (2.6pt); }
\draw[structurepurple, very thick, -{Stealth[length=2.5mm]}] (4,3) -- ({4+1.342},{3+1.342});
\draw[targetgreen!70!black, very thick, -{Stealth[length=2.5mm]}] (4,3) -- ({4+0.447},{3-0.447});
\node[structurepurple, font=\small, anchor=west, fill=white, inner sep=1pt] at (5.45,4.5) {PC1: variance $3.6$};
\node[targetgreen!70!black, font=\small, anchor=north west, fill=white, inner sep=1pt] at (4.55,2.5) {PC2: variance $0.4$};
\node[font=\small, anchor=east] at (3.55,2.4) {mean};
\draw[gray] (3.55,2.5) -- (3.94,2.94);
\node[densityred, font=\small, anchor=north west] at (5.2,1.1) {2-SD Normal contour};
\fill (4,3) circle (1.5pt);
\end{tikzpicture}
\caption{Six points with covariance matrix $S$. The arrows are the eigenvectors of $S$, each
drawn from the mean $(4,3)$ with length equal to its standard deviation ($\sqrt{3.6}\approx1.90$ and
$\sqrt{0.4}\approx0.63$). The red ellipse is the 2-SD contour of a Normal with this
covariance; its long axis is PC1. PCA keeps that axis and drops the short one.}
\label{fig:ml-pca}
\end{figure}

\begin{connectionbox}
PCA, Gaussian mixture models, Gaussian processes, and the Kalman filter are all built on
covariance matrices. The initialisation schemes and normalisation layers in deep learning are
attempts to keep the covariance of activations close to the identity matrix. ISLR Ch.~12 and
Bishop Ch.~12 cover PCA in full.
\end{connectionbox}

\section{What you can now read}
\label{sec:ml-reading}

The statistical half of the standard ML texts is now familiar ground. Table~\ref{tab:ml-map}
maps their early chapters to this book.

\begin{table}[H]
\centering
\small
\renewcommand{\arraystretch}{1.25}
\begin{tabular}{p{4.6cm} p{4.6cm} p{3.6cm}}
\toprule
\textbf{Text and chapter} & \textbf{Topic} & \textbf{Where it is here} \\
\midrule
ISLR Ch.~2 & bias--variance, train vs test MSE & Ch.~14, Sec.~\ref{sec:ml-biasvar} \\
ISLR Ch.~3--4 & linear and logistic regression & Ch.~12--13 \\
ISLR Ch.~5 & cross-validation, bootstrap & Ch.~8, Sec.~\ref{sec:ml-splits} \\
ISLR Ch.~6 & ridge, lasso, model selection & Sec.~\ref{sec:ml-reg} \\
ISLR Ch.~12 & PCA & Sec.~\ref{sec:ml-cov} \\
Bishop Ch.~1 & curve fitting, overfitting, Bayes, decision theory & Ch.~14, Secs.~\ref{sec:ml-risk}--\ref{sec:ml-biasvar} \\
Bishop Ch.~2 & Bernoulli, Binomial, Gaussian, MLE & Ch.~6, 14, Sec.~\ref{sec:ml-cov} \\
Bishop Ch.~3 & linear basis models, regularisation, bias--variance & Ch.~12, Secs.~\ref{sec:ml-biasvar}--\ref{sec:ml-reg} \\
Goodfellow Ch.~3 & probability, expectation, covariance, Gaussians & Sec.~\ref{sec:ml-notation}, \ref{sec:ml-cov} \\
Goodfellow Ch.~5 & capacity, validation, estimators, bias/variance, MLE, SGD & this whole chapter, Ch.~14 \\
\bottomrule
\end{tabular}
\caption{Reading map. ISLR is James, Witten, Hastie and Tibshirani,
\emph{An Introduction to Statistical Learning} (2nd ed.). Bishop is \emph{Pattern Recognition
and Machine Learning}. Goodfellow is Goodfellow, Bengio and Courville, \emph{Deep Learning}.}
\label{tab:ml-map}
\end{table}

When a page gets dense, apply the decoder: find each expectation and ask ``average of what, over
which distribution?''; find each hat and ask ``estimated from which data?''; find each reported
number and ask ``what is its standard error?''

\section{Summary}

\begin{itemize}
  \item $\E_{x\sim p}$ is an average over the world; $\frac1n\sum_i$ is the same average over
  your data ($\E_{x\sim\hat p}$). Hats mark estimates; $\argmin$ is a location.
  \item ERM minimises empirical risk. Its training loss is optimistic; the generalisation gap
  $R(\hat f)-\hat R_n(\hat f)$ grows with flexibility.
  \item Train / validation / test: fit, choose, report. $k$-fold CV reuses data for choosing.
  Every data-dependent preprocessing step goes inside the split.
  \item Expected test error $=\sigma^2 + \text{bias}^2 + \text{variance}$. Flexibility trades
  bias for variance and the sum is U-shaped.
  \item Ridge divides coefficients toward zero; lasso subtracts toward zero and produces exact
  zeros (corners of the diamond).
  \item From the confusion matrix: precision $=TP/(TP+FP)$, recall $=TP/(TP+FN)$,
  $F_1$ = harmonic mean. AUC $=\P(\text{score}_+>\text{score}_-)$. Precision depends on the
  base rate; ROC does not. Calibration is a separate property, read off a reliability diagram.
  \item Metrics have standard errors: binomial CI for accuracy, bootstrap for AUC, McNemar or
  paired bootstrap to compare two models on one test set.
  \item A minibatch gradient is an unbiased Monte Carlo estimate of the full gradient with SE
  $\propto 1/\sqrt B$.
  \item The covariance matrix holds variances and covariances; its eigenvectors are the
  principal components and the axes of multivariate Normal contours.
\end{itemize}

Chapter 16 collects the formulas from the whole book in one place, together with a decision
guide for choosing an interval or a test.

\section{Exercises}

\begin{enumerate}
  \item A fraud model on $1000$ transactions has $TP=45$, $FP=15$, $FN=5$, $TN=935$. Compute
  accuracy, precision, recall and $F_1$. What accuracy does ``never flag'' get?
  \item A classifier gets $368$ of $400$ test examples right. Give a $95\%$ confidence interval
  for its accuracy.
  \item Models A and B are scored on the same test set. A is right and B wrong on $30$ examples;
  B is right and A wrong on $12$. Compute McNemar's $\chi^2$ and its $p$-value.
  \item Two positives score $0.9$ and $0.4$; two negatives score $0.6$ and $0.2$. Compute the
  AUC by counting pairs, and draw the ROC curve.
  \item With an orthonormal design, least squares gives $\hat\beta=(2.0,\,-0.3,\,0.8)$. Find
  the ridge and lasso estimates for $\lambda=1$ using the formulas of
  Section~\ref{sec:ml-reg}.
  \item A screening test has $\text{TPR}=0.9$ and $\text{FPR}=0.05$. The condition affects $2\%$
  of people tested. What fraction of positive results are true positives?
  \item A dataset has covariance matrix $\begin{pmatrix}4&0\\0&1\end{pmatrix}$. What are the
  principal components, and what fraction of the variance does the first one capture?
  \item Show that a minibatch gradient computed from $B$ examples drawn uniformly at random with
  replacement is unbiased for the full-batch gradient, and that its variance is $\sigma^2/B$
  where $\sigma^2$ is the variance of the per-example gradients.
  \item A team standardises all features using the full dataset's mean and SD, then runs
  $5$-fold CV. Another team picks the $10$ features most correlated with the label on the full
  dataset, then runs $5$-fold CV. Which leak is worse, and why?
  \item How many test examples do you need for a $95\%$ CI of half-width $0.02$ on an accuracy
  near $0.8$?
\end{enumerate}

\subsection*{Answers}

\begin{enumerate}
  \item Accuracy is $980/1000=0.98$, precision is $45/60=0.75$, and recall is
  $45/50=0.90$. So $F_1=1.35/1.65\approx0.818$. Never flagging scores $0.95$.
  \item $\hat a=0.92$, $\SE=\sqrt{0.92\times0.08/400}\approx0.0136$,
  interval $0.92\pm0.027=[0.893,\ 0.947]$.
  \item $\chi^2=(30-12)^2/42=324/42\approx7.71$, $p\approx0.005$. Strong evidence that A is
  better on this distribution.
  \item Pairs: $0.9>0.6$, $0.9>0.2$, $0.4>0.2$, and $0.4<0.6$. AUC $=3/4=0.75$. ROC path:
  $(0,0)\to(0,0.5)\to(0.5,0.5)\to(0.5,1)\to(1,1)$.
  \item Ridge: $(1.0,\,-0.15,\,0.4)$. Lasso: $(1.5,\,0,\,0.3)$.
  \item $\dfrac{0.9(0.02)}{0.9(0.02)+0.05(0.98)}=\dfrac{0.018}{0.067}\approx0.27$. About
  three positives in four are false.
  \item The coordinate axes: PC1 along $x_1$ (variance $4$), PC2 along $x_2$ (variance $1$).
  PC1 captures $4/5=80\%$.
  \item Each sampled index $I_j$ is uniform on $\{1,\dots,n\}$, so
  $\E[g_{I_j}]=\frac1n\sum_i g_i=\bar g$, and the batch mean has expectation $\bar g$ by
  linearity. The $B$ draws are independent, each with variance $\sigma^2$, so the variance of
  their mean is $\sigma^2/B$ (Chapter 6).
  \item Feature selection is far worse. Standardising leaks only a mean and SD per feature,
  which barely changes. Selecting features by their correlation with the label lets every fold's
  labels influence which features the model sees, which can produce high CV accuracy on pure
  noise.
  \item $n=1.96^2(0.8)(0.2)/0.02^2\approx1536.6$, so about $1540$ examples.
\end{enumerate}
